\section{Problem Formulation}
\label{sec:problem}

In this section, we formally define the problem of parameter estimation for dynamical systems under \emph{partial and sparse observation} settings. 
Let $\mathbf{x}(t)\in\mathbb{R}^d$ denote the full state of the dynamical system and $\theta\in\Theta\subset\mathbb{R}^p$ the unknown parameter vector.
The underlying dynamics are governed by the ordinary differential equation (ODE)
\begin{equation}
\label{eq:ode}
    \dot{\mathbf{x}}(t)=f(\mathbf{x}(t),\theta),\qquad t\in[0,T],
\end{equation}
where $f:\mathbb{R}^d\times\Theta\to\mathbb{R}^d$ is a known, sufficiently smooth vector field.
The classical parameter estimation problem seeks to infer the unknown parameter $\theta$ given the continuous trajectory of the full state $\mathbf{x}(t)$~\cite{tarantola2005inverse}. 
In practical applications, however, obtaining such complete observations is rarely feasible.
The following subsections define the specific constraints of realistic observations and formulates the resulting inverse problem.

\subsection{Partial and Sparse Observation Setting}
\label{subsec:obsmodel}

We consider a \emph{partial observation} setting, in which only a subset of the state variables is observable.
We decompose the full state $\mathbf{x}(t)$ into an \emph{observed} state $x_{\mathrm{obs}}(t)\in\mathbb{R}^{d_{\mathrm{obs}}}$ and a
\emph{unobserved} state $x_{\mathrm{hid}}(t)\in\mathbb{R}^{d_{\mathrm{hid}}}$, such that $d_{\mathrm{obs}}+d_{\mathrm{hid}}=d$.
We refer to this unobserved state as the \emph{hidden state}. Reconstructing it explicitly is a structural choice --- an inductive bias --- of the decoupled estimator we study, rather than a requirement for inferring the unknown parameter $\theta$; a direct estimator infers $\theta$ without forming it, and Section~\ref{sec:correctness} examines when this reconstruction actually helps.
Without loss of generality, the state vector can be partitioned as $\mathbf{x}(t) = \left[x_\mathrm{obs}(t)^\top , x_\mathrm{hid}(t)^\top\right]^\top$.
To formalize this limitation, we define a non-injective map $h:\mathbb{R}^d \to \mathbb{R}^{d_{obs}}$ acting as a coordinate projection:

\begin{equation}
\label{eq:h-projection}
    x_{\mathrm{obs}}(t)=h(\mathbf{x}(t)).
\end{equation}



The data are further constrained by \emph{sparse sampling} at discrete times.
We assume a uniform sampling scheme over the time interval $[0, T]$ with a fixed step size $\Delta t$:

\begin{equation}
\label{eq:sampling-times}
    0=t_0<t_1<\cdots<t_{N}=T,
\end{equation}

where $N\Delta t = T$.
The term \emph{sparse} specifically denotes scenarios where $(N+1)$ is extremely small relative to the intrinsic time scale of the dynamics.
Consequently, the partial and sparse time-series data collected from the system is denoted as the set of discrete observations $\left\{x_\mathrm{obs}(t_i)\right\}_{i=0}^{N}$.

Since our estimators are modeled as multi-layer perceptrons (MLPs), the discrete sampled trajectories must be flattened into 1D vectors.
We write the flattened observation vector $\mathbf{x}_\mathrm{obs} \in \mathbb{R}^{(N+1)
\times d_\mathrm{obs}}$ and the hidden state vector $\mathbf{x}_\mathrm{hid} \in \mathbb{R}^{(N+1) \times d_\mathrm{hid}}$ by concatenating the state vectors at each sampling time:

\begin{align}
    \mathbf{x}_{\mathrm{obs}}&
    :=[x_{\mathrm{obs}}(t_0)^\top,\ldots,x_{\mathrm{obs}}(t_{N})^\top]^\top
,
\\
    \mathbf{x}_{\mathrm{hid}}&
    :=[x_{\mathrm{hid}}(t_0)^\top,\ldots,x_{\mathrm{hid}}(t_{N})^\top]^\top.
\end{align}

Estimating the unknown parameter vector $\theta$ from the partial and sparse observations $\mathbf{x}_\mathrm{obs}$ constitutes an inverse problem involving unobserved state trajectories. 
Since the hidden state $\mathbf{x}_\mathrm{hid}$ is unmeasured yet continuously coupled with the observed dynamics through the governing ODE, traditional joint optimization over both states and parameters frequently leads to highly non-convex and ill-conditioned objective landscapes. 
This computational difficulty is significantly amplified in sparse sampling regimes, where classical solvers are highly sensitive to initialization and prone to getting trapped in local minima or diverging. 
To alleviate these numerical challenges, a structured inference framework that systematically decouples the state reconstruction and parameter estimation tasks is required.

\begin{remark}[Why sparsity is essential]
Given the \emph{continuous} partial observation $x_\mathrm{obs}(t)$ for all $t\in[0,T]$, classical theory offers
recovery guarantees: Takens' embedding theorem~\cite{takens1981detecting} and the Hermann--Krener observability rank condition~\cite{hermann1977nonlinear}
(Section~\ref{sec:identifiability}) certify, under genericity assumptions, that the full state and the
parameters can in principle be reconstructed from the observed component and its time derivatives.
These guarantees rely on access to high-order temporal information (derivatives, or equivalently a dense
time series). Under the sparse sampling of \eqref{eq:sampling-times}, that information is unavailable,
since derivatives cannot be reliably estimated from a handful of points. The guarantees then no longer
apply, and the problem may become practically, or even structurally, non-identifiable. This is the
regime we study, and it is what makes the inverse problem genuinely hard.
\end{remark}